\section{Publication, compression and persistent state}
\label{architecture-publication}

\subsection{A frame is published by a protocol}

The output layer translates a completed model state into selected products.
It resolves the field inventory and units, stages required values, encodes
their storage representation and publishes metadata identifying the completed
frame. The existence of a directory or one field file is not enough to tell
a reader that all required fields belong to one finished step. Studio's frame
discovery uses publication state and manifest identity to avoid displaying a
mixture of old and partially written new data.

\href{../src/core/output/output_writer.cpp}{AsyncOutputWriter} has synchronous
and asynchronous paths. A synchronous writer can use a valid borrowed view
while the caller holds the state stable. A background writer requires owned
snapshot storage; \texttt{submit} rejects a borrowed field rather than letting
the integration thread mutate data while compression reads it. Its queue has
a one-frame ownership bound, and the ordinary runtime waits for availability
before assembling the next full snapshot. Backpressure is deliberate: an
unbounded queue would convert a slow filesystem into uncontrolled RAM growth.

An asynchronous writer does not remove serialization or compression work. It
allows some work to overlap while introducing a snapshot, queue synchronization
and a separate error path. A write failure sets an error state and wakes waiting
participants so that an output error cannot strand the producer indefinitely.
The scientifically relevant frame remains identified by model time and source
state, not by the wall-clock moment when the filesystem completes its write.

NPY provides a direct array format, netCDF provides structured scientific
history, and ZFP provides selected compressed representations. Their wrappers
must preserve axis order, shape, field identity and physical metadata. Internal
storage order need not be the order a downstream VTK image expects; that
conversion belongs to the reader/preparation boundary. A decoder that produces
the right number of floats in the wrong order has not preserved a field.

\subsection{What ZFP contributes, and what Nimbus adds}

Peter Lindstrom's 2014 paper, \emph{Fixed-Rate Compressed Floating-Point Arrays},
develops a block representation intended to make compressed numerical arrays
accessible at a prescribed storage rate. Small $4^d$ blocks in $d$ dimensions
share a floating-point representation, undergo a decorrelating transform and
are encoded by significance. The paper supplies the codec method; it does not
specify Nimbus's history container, keyframe schedule, field tolerances or
publication transaction. See the
\href{https://vis.cs.ucdavis.edu/vis2014papers/TVCG/papers/2674_20tvcg12-lindstrom-2346458.pdf}{primary paper}.

The current library's modes answer different questions. Fixed rate constrains
bits per value/block, while error depends on the data. Fixed precision controls
encoded transform precision; it is not a uniform pointwise relative-error
promise for every value in a mixed-magnitude block. Fixed accuracy targets an
absolute error scale in the field's units and produces variable-length output.
Reversible mode supplies bit-preserving reconstruction. These mode distinctions
are defined by the versioned
\href{https://zfp.readthedocs.io/en/release1.0.1/modes.html}{ZFP 1.0.1 mode reference}.
The library also documents finite-representation limits for exceptionally small
accuracy tolerances. Zero tolerance in a lossy mode must not be described as
equivalent to reversible compression; see
\href{https://zfp.readthedocs.io/en/release1.0.1/faq.html}{FAQ questions 17, 20 and 21}.

Nimbus calls the corresponding stream setters in its output writer and records
the mode parameters in its own container. Ordinary public mode parsing accepts
\texttt{accuracy}, \texttt{precision} and \texttt{rate}; admission imposes
additional restrictions for a supported scientific profile. The retained
ordinary admitted lossy-history contract requires accuracy mode and per-field
tolerances, prohibits a positive sparse threshold and disallows the optional
Float16 prefilter. Configuration code also caps the potential-vorticity
tolerance at $10^{-8}$ in its SI representation: a generic visual tolerance
could erase a meteorologically significant signal while still being a valid
codec setting. These are Nimbus decisions in
\href{../src/core/output/output_config.cpp}{output configuration}, not universal
recommendations from the compression paper.

The optional sparse and Float16 transforms are separate lossy operations before
compression. The former sets eligible small values to zero; the latter rounds
eligible values to reduced significand precision. A codec tolerance applies to
the values supplied to that codec. It does not retroactively bound information
already removed by a prefilter. For a source $x$, prefilter $F$, codec
reconstruction $D(C(F(x)))$ and later reconstruction $R$, a useful accounting
separation is
\[
 |x-\widehat x|
 \leq E_{\mathrm{prefilter}}+E_{\mathrm{codec}}+E_{\mathrm{reconstruction}},
\]
when the individual terms have actually been defined and bounded for that
path. Writing down the inequality does not prove all three bounds. This is why
the admitted policy restricts transforms and why quantitative output checks
must compare the final decoded field, not just inspect a mode name.

Temporal delta history is another wrapper-level choice. The encoder keeps
decoded previous fields so that the reader and writer predict from the same
represented history. For a linear two-frame predictor, the conceptual equations
are
\[
\begin{aligned}
 P_n &= 2\widehat x_{n-1}-\widehat x_{n-2},\\
 r_n &= x_n-P_n,\\
 \widehat x_n &= P_n+\widehat r_n.
\end{aligned}
\]
The implementation computes residual and reconstruction expressions with
specified wider intermediates and Float32 storage. With identical predictor
history, old error is already incorporated in the residual definition; using
uncompressed previous fields at the writer and decoded previous fields at the
reader would introduce a different drift mechanism. Actual error still
includes residual rounding, codec error and final reconstruction rounding.
Keyframes reset history dependencies, and retention cannot remove a frame
that a retained delta chain still needs.

Nimbus's spatial container bricks are also different from ZFP's internal
$4^d$ codec blocks. The current bricked reader can select $64\times64\times32$
container bricks, inspect a directory of offsets and reconstruct only the
required spatial region for supported versions. Legacy containers can require
full-volume reconstruction. A small requested output window therefore does
not always imply a small decode working set. The native reader reports these
capabilities instead of silently labeling a full decode as a cheap window.

Field-parallel publication assigns independent fields to workers with isolated
codec workspaces; libzfp execution within each field remains serial in this
policy. It is not the same as enabling the library's internal parallel codec.
Scratch budgets and exact serial/parallel artifact checks are part of the
wrapper contract. More workers may improve throughput, increase memory
pressure, or do neither; the retained benchmark harness makes no universal
throughput claim. The
\href{../tests/core/test_zfp_roundtrip.cpp}{roundtrip tests} and
\href{../tests/integration/test_zfp_field_parallel_benchmark.py}{benchmark
contract tests} address different parts of this behavior.

The reversible radar-v4 container uses the library's reversible mode and
per-brick integrity information under a separate guarded mixed-history path.
It must not be used to claim that ordinary lossy output is exact, or that the
general radar-v4 readiness work is complete. Output history is also not the
solver's in-memory representation: this architecture does not generally
compress prognostic arrays in place between numerical operators. Paper results
about compressed numerical-array computation cannot be transferred to Nimbus's
history writer without a separate implementation and measurement.

\subsection{Checkpoints encode continuation, not just visible fields}

A restart needs more than wind, temperature and water arrays. It can require
scheme-private particle state, time/cadence counters, boundary history, nested
execution state, trajectories and other runtime owners. The checkpoint layer
serializes a versioned set of sections rather than writing the raw bytes of a
C++ object graph. Raw pointers, allocator state, vtable addresses and mutexes
have no portable meaning in another process.

\href{../include/core/runtime/checkpoint_format.hpp}{The wire format} uses
explicit little-endian integer encodings and IEEE floating-point values,
section identifiers, versions, lengths, flags and checksums. It rejects
duplicate or unsupported required sections and checks sizes before allocation.
The architecture-specific in-memory layout is therefore separated from the
continuation format. A native reader can convert byte order where needed;
compatibility still depends on supported versions and the semantic state
contract, not merely on recognizing the file magic.

The writer streams sections into a temporary file while computing integrity
information, closes and synchronizes it, atomically publishes the destination,
then synchronizes the containing directory through platform helpers. This
avoids treating a partial file as the final checkpoint. On restore, decoded
candidate state and attachments are validated before replacing the live state.
An unknown future section or inconsistent tracking authority must not become
a partially restored simulation. See
\href{../src/core/runtime/simulation_checkpoint.cpp}{checkpoint assembly and restore},
\href{../tests/core/test_checkpoint_wire.cpp}{wire tests} and
\href{../tests/core/test_simulation_checkpoint.cpp}{continuation tests}.

\section{Studio IPC, asynchronous work and remote execution}
\label{architecture-ipc}

\subsection{Keep native resources behind the Rust boundary}

Studio's React components express user intent and maintain view state. Rust
owns native filesystem and process operations. Tauri commands carry typed
requests and responses across that boundary. A field-window response uses a
compact binary envelope with a small validated header and Float32 payload,
rather than inflating millions of values into a JSON number array. Geometry
may add aligned Float64 components when required by the declared schema.

The native command checks source dimensions, range, stride, output count,
arithmetic overflow and byte limits before serving the request. Its current
window limits distinguish 64 MiB of returned field data from a larger bounded
source working set. A shared native I/O reservation budgets source bytes plus
twice the output bytes under a 512 MiB aggregate ceiling, then releases the
claim through RAII when the last owner drops it. This models simultaneous
work, not just the size of an individual successful response. See
\href{../ui/src-tauri/src/commands/outputs.rs}{native field commands} and
\href{../ui/src-tauri/src/commands/io_budget.rs}{I/O reservations}.

These values are implementation ceilings, not a promise that every computer
can safely devote that amount to Studio. Frontend preparation and GPU textures
have additional budgets. A renderer can be within its texture-size limit yet
run out of combined native, JavaScript and GPU working memory. The system must
retain the reason a representation was reduced or refused rather than silently
labeling it full resolution.

\subsection{Cancellation is a state transition, not instantaneous preemption}

Each native window read has a request identity, cancellation flag and phase.
Cooperative loops check the flag during header reading, streaming and other
interruptible work. A libzfp decode or native netCDF hyperslab call cannot be
interrupted at an arbitrary instruction through this interface. Cancellation
is checked at its boundaries; a cancelled result is then discarded. The
receipt says cancellation was requested, not that an in-progress native call
has already stopped.

The frontend has a complementary obligation: a late result from frame $n$ must
not replace frame $n+1$ after the user scrubs forward. Request/generation
identity and coherent scene publication solve a different problem from the
native cancellation flag. Even a request too late to cancel can be prevented
from publishing into the wrong view. The
\href{../ui/src/lib/run/frameIdentity.ts}{frame identity} and
\href{../ui/src/lib/vtk/coherentSceneGeneration.ts}{scene generation}
owners make that distinction explicit.

JavaScript workers perform preparation and extraction outside the main UI
thread. Transferable ArrayBuffers move ownership and detach the transferred
buffer; SharedArrayBuffers share storage when the environment permits them.
Nimbus's handoff protocol distinguishes a transferred render-only result from
a retained scientific source and from a shared immutable dataset. The choice
must preserve retries and cache users. Avoiding a structured-clone copy is
useful, but does not establish that native IPC or GPU texture upload was
zero-copy.

\subsection{Launching and stopping a process}

The \href{../ui/src-tauri/src/commands/runner.rs}{native process runner}
checks the executable and configuration identities before spawning. It
publishes immutable launch provenance and starts separate readers for stdout
and stderr plus a waiter for process completion. This keeps a child blocked
on a full output pipe from appearing to be a stalled numerical solver. The
process exit state is distinct from a complete final frame or a verified
checkpoint; readers must still inspect publication records.

Stopping is also owned. On applicable Unix paths, the runner requests graceful
termination and escalates if the child does not exit within its grace period.
It does not equate closing a plot with killing the simulation, and a failed
process does not become a successful run merely because it emitted several
valid early frames. Filesystem, process and application lifecycle operations
must preserve their separate error states.

Remote execution stages an exact request on a user-configured SSH target,
checks the remote executable's catalog and identities, and launches that
process there. Receiving output has its own inventory, per-file hashes and
completion/cancellation receipt. A cancelled transfer is not a complete local
copy of the remote run. The
\href{../ui/src-tauri/src/commands/remote_compute.rs}{remote compute owner}
and \href{../ui/src-tauri/src/commands/remote_receive.rs}{receive owner}
separate launching from verified retrieval.

This is remote process orchestration, not MPI domain decomposition. The current
architecture does not automatically distribute one atmospheric grid across
multiple networked hosts. SSH availability, remote CPU ISA, installed library
ABI, GPU capabilities and local viewer support remain separate checks. The
workstation being unavailable prevents new measurements there; it does not
justify inferring their results from a successful local ARM build.
